\subsection{Công Cụ Trực Quan}
% -----------------------------------------------------------------------------

\subsubsection{Biểu Đồ Tần Suất của Phần Dư}

Dấu hiệu vi phạm: \textbf{lệch phải/trái} (đuôi một bên dài), \textbf{đuôi nặng} (đỉnh nhọn, đuôi dày), \textbf{đa đỉnh} (gợi ý dữ liệu từ 2 quần thể).

\subsubsection{Q-Q Plot (Quantile-Quantile Plot)}

\vspace{4pt}
\begin{center}
\begin{tabular}{@{}ll@{}}
\toprule
\textbf{Mẫu hình quan sát} & \textbf{Diễn giải} \\
\midrule
Các điểm trên đường thẳng & Phân phối chuẩn \\
Hai đầu lệch lên trên (dạng chữ S) & Đuôi nặng hai phía \\
Hai đầu lệch xuống dưới & Đuôi mỏng \\
Đuôi phải lệch lên & Lệch phải \\
Đuôi trái lệch xuống & Lệch trái \\
\bottomrule
\end{tabular}
\end{center}

\textbf{Lưu ý.} Đồ thị xác suất chuẩn nhạy hơn biểu đồ tần suất trong việc phát hiện lệch ở đuôi phân phối; ưu tiên dùng khi mẫu nhỏ.

% -----------------------------------------------------------------------------
